% Generated draft: requires mathematical and source review.
\documentclass[11pt]{article}
\usepackage[margin=25mm]{geometry}
\usepackage{fontspec}
\setmainfont{Latin Modern Roman}
\usepackage{amsmath,amssymb,mathtools,unicode-math}
\setmathfont{Latin Modern Math}
\usepackage{enumitem,needspace,xurl,hyperref}
\setlength{\parindent}{0pt}
\setlength{\parskip}{5pt}
\setlength{\emergencystretch}{3em}
\newcommand{\N}{\mathbb{N}}
\newcommand{\Z}{\mathbb{Z}}
\newcommand{\Q}{\mathbb{Q}}
\newcommand{\R}{\mathbb{R}}
\newcommand{\C}{\mathbb{C}}
\newcommand{\F}{\mathbb{F}}
\newcommand{\A}{\mathbb{A}}
\newcommand{\PP}{\mathbb{P}}
\newcommand{\E}{\mathbb{E}}
\newcommand{\eps}{\varepsilon}
\newcommand{\dd}{\,\mathrm{d}}
\begin{document}
\section*{UnsolvedMath ID: 30003956}
\textbf{Source problem code:} catalogue entry
\subsection{Definitions and mathematical statement}
Let \(S\) denote the underlying surface, \(\mathcal T(S)\) its Teichm\"uller space of marked conformal structures, and \(\operatorname{Mod}(S)\) its mapping class group, acting by changing markings. The supplied record does not specify the topological type of \(S\), boundary conditions, or puncture conventions. The following formulation is conditional on choosing these data so that the standard hyperbolic-length definition of the Thurston metric below applies; it does not resolve those missing hypotheses.

Let \(\mathcal C(S)\) be the set of free homotopy classes of essential, nonperipheral simple closed curves on \(S\). For \(z\in\mathcal T(S)\) and \(\alpha\in\mathcal C(S)\), write \(\ell_z(\alpha)\) for the length of the geodesic representative in the hyperbolic structure associated with \(z\). In the usual setting where this formula defines the Thurston metric, put
\[
 d_{\mathrm{Th}}(u,v)
 =\log\sup_{\alpha\in\mathcal C(S)}
 \frac{\ell_v(\alpha)}{\ell_u(\alpha)}.
\]
This metric is generally asymmetric. Consequently, the phrase ``ball about \(x\)'' requires a direction convention. For closed balls, the two corresponding counting functions are
\[
\begin{aligned}
 N^+_{x,y}(R)
 &=\#\{z\in\operatorname{Mod}(S)\cdot y:
             d_{\mathrm{Th}}(x,z)\le R\},\\
 N^-_{x,y}(R)
 &=\#\{z\in\operatorname{Mod}(S)\cdot y:
             d_{\mathrm{Th}}(z,x)\le R\},
 \qquad R\ge0.
\end{aligned}
\]
These count distinct orbit points, with no multiplicity from the stabilizer of \(y\). In particular, the forward condition is
\[
 \ell_z(\alpha)\le e^R\ell_x(\alpha)
 \quad\text{for every }\alpha\in\mathcal C(S).
\]
The reverse condition instead requires
\(\ell_x(\alpha)\le e^R\ell_z(\alpha)\) for every such \(\alpha\).

\textbf{Requested asymptotic problem.} For arbitrary fixed \(x,y\in\mathcal T(S)\), determine the asymptotic behavior, as \(R\to\infty\), of the orbit count for the intended Thurston ball. The supplied wording does not select \(N^+_{x,y}\) or \(N^-_{x,y}\), and does not specify open versus closed balls. The closed-ball formulas above make the directional alternatives explicit without asserting that the source requests both.

For comparison, the extremal length of \(\alpha\) on the conformal surface \(z\) is
\[
 \operatorname{Ext}_z(\alpha)
 =\sup_{\rho}
 \frac{L_\rho(\alpha)^2}{\operatorname{Area}_\rho(z)},
\]
where \(\rho\) ranges over conformal metrics of positive finite area and \(L_\rho(\alpha)\) is the infimum of the \(\rho\)-lengths of curves representing \(\alpha\). The source further asks how to relate or unify the Thurston orbit asymptotic with counting by extremal length rather than hyperbolic length. It does not specify the objects counted in that comparison, a chosen curve or collection of curves, the extremal-length functional, or its threshold normalization. Thus this comparison remains an unresolved part of the formulation, rather than a uniquely defined second counting problem.
\subsection{Short English statement}
For fixed points \(x,y\) in a Teichm\"uller space, determine how the number of distinct mapping-class orbit points of \(y\) in a Thurston ball about \(x\) grows as its radius tends to infinity, and relate this to extremal-length counting. The surface conventions, direction of the asymmetric ball, and precise extremal-length comparison require clarification.
\subsection{Sources}
\begin{itemize}
\item Local registry: \texttt{data/unsolvedmath/problems.json}, with supplementary records in \texttt{catalog/problems.json}; supplied permanent UnsolvedMath ID \(30003956\), source problem code \texttt{catalogue entry}, title ``Mapping-Class Orbit Counting in the Thurston Metric.''
\item Supplied report citation: \emph{New Trends in Teichm\"uller Theory and Mapping Class Groups} (2019), \emph{Oberwolfach Reports}, pp.\ 2475--2534, DOI \texttt{10.4171/owr/2018/40}: \url{https://doi.org/10.4171/owr/2018/40}. The record separately gives proposed year 2018; these year fields are retained without reconciling them.
\item Supplied literature source labeled ``Oberwolfach Report 40, Su problem'': \url{https://publications.mfo.de/bitstream/handle/mfo/3662/OWR_2018_40.pdf?isAllowed=y&sequence=1}. This label associates the problem with Su; the record's proposer field is empty.
\item Supplied literature source labeled ``Su, variation of extremal length'': \url{https://arxiv.org/abs/1210.0743}.
\item Supplied literature source labeled ``Arana-Herrera, effective mapping class dynamics II'': \url{https://link.springer.com/article/10.1007/s00039-026-00740-6}.
\item Supplied literature assessment, checked \texttt{2026-08-21}: the record labels the problem open and reports known Teichm\"uller-metric orbit counting, partial Thurston-metric and extremal-length results, and unproved requested extremal-length asymptotics and unification. This is supplied dated curation derived from the OWR record, not an independently verified literature assessment. The supplied background explicitly requests independent verification of time-sensitive status and cited claims.
\end{itemize}
\end{document}
